% ============================================================
%  lang/en.tex --- every user-visible string in the English edition
%
%  Nothing in chapters/en/ may hard-code a word that a macro also typesets.
%  If a label is needed and is not here, add it here AND in lang/pl.tex; CI
%  fails the build when the two files do not define the same macro set (C3).
% ============================================================

% ---------- Part titles ----------
% Part titles live here rather than in structure.tex: nothing outside a
% language file may hard-code a word the other edition also sets.
\newcommand{\lblPartI}{Runtime and toolchain}
\newcommand{\lblPartII}{The language, mapped}
\newcommand{\lblPartIII}{Concurrency}
\newcommand{\lblPartIV}{Shipping}
\newcommand{\lblPartV}{Python for AI work}
\newcommand{\lblPartAppendices}{Appendices}

% ---------- Admonition titles ----------
\newcommand{\lblNote}{Note}
\newcommand{\lblWarning}{Watch out}
\newcommand{\lblTrap}{The trap}
\newcommand{\lblCsharp}{In C\# terms}
\newcommand{\lblVersion}{Version sensitive}
\newcommand{\lblVerify}{Run this before you trust it}
\newcommand{\lblExercise}{Exercise}
\newcommand{\lblProject}{trace-assert stage}
\newcommand{\lblMeasured}{Measured}

% ---------- The exercise footer ----------
% An exercise names the file the reader opens in the editor beside the PDF,
% the test that decides whether it is done, and the one command that runs it.
\newcommand{\lblExerciseStarter}{Starter}
\newcommand{\lblExerciseTest}{Test}
\newcommand{\lblExerciseRun}{Run}

% ---------- Build-state markers ----------
\newcommand{\lblNotWritten}{NOT YET WRITTEN}
\newcommand{\lblNotRendered}{NOT RENDERED --- run \texttt{make diagrams}}
\newcommand{\lblSourceMissing}{Mermaid source missing}
\newcommand{\lblTranscriptMissing}{TRANSCRIPT NOT GENERATED --- run \texttt{make numbers}}
\newcommand{\lblValueMissing}{value missing}
\newcommand{\lblDiagramManifest}{Diagram manifest}
\newcommand{\lblExerciseManifest}{Exercise manifest}

% ---------- Front-matter running heads ----------
\newcommand{\lblHowToUse}{How to use this book}
\newcommand{\lblIntroduction}{Introduction}

% The punctuation dash. English sets an em dash closed up; Polish sets a
% spaced en dash. Authors write \dash and never type a dash directly.
\newcommand{\dash}{---}

% ---------- Index cross-references ----------
\newcommand{\lblIndexSee}{see}
\newcommand{\lblIndexSeeAlso}{see also}

% The date the pins were last checked against PyPI. User-visible, so it lives
% here rather than in preamble.tex.
\newcommand{\pinnedon}{September 2026}

% ---------- PDF metadata ----------
% The subject, keyword list and document language differ between the editions.
% The PDF *title* stays in the main files, because those are the wiring.
\newcommand{\lblPdfSubject}{Reading, writing and shipping Python when C\# is your first language}
\newcommand{\lblPdfKeywords}{python, dotnet, csharp, career transition, uv, pydantic, fastapi, sqlalchemy, pytest, asyncio}
\newcommand{\lblPdfLang}{en-GB}
